\chapter{章程总目}
\chapterkana{章程総目}
\epigraphbox{“萧规曹随。”}{——《史记·曹相国世家》}

\kaishou{一本管了两万年的章程，每一页的末尾都留着一栏，用来记改动。\par
一千二百四十七条，那一栏全空。}

\section{四百一十二枚}

一九六五年，苏联的一次无人月球采样带回了一批物件，装在一只木箱中。箱面登记得相当含糊，只写着「远东出土文书」，来源不明，暂不归类。那几年正值美苏轮番向月面发射探测器，带回的样品一批接一批，事后谁也说不清这一箱究竟出自哪一次任务。七十年代这批材料并入基金会编目时，木箱原封未动，登记也照旧抄录。

箱内盛放着四百一十二枚金属片，均为单面刻字，最长者十八厘米，最宽者九厘米，厚度不足一毫米。每枚边缘都有穿孔，孔距不等，孔沿磨得发亮，原应是以丝线串连成册，如今丝线早已朽尽，只剩穿孔。\footnote{金属片的成分、锈蚀分布与库房编号，见基金会的移交清单；本书只用得着其中三项：片数、尺寸、穿孔。}随箱附纸一页，上书「顺序不能确定，暂依厚度排列」。编目者写下这句话时透着无奈：一部成册的典籍散成了四百一十二片零金碎玉，谁也分不清哪一枚该接哪一枚。今日通行的卷次与条序，其实是后人按内容重新排定的。

重排之后，四百一十二枚金属片分作六卷，合计一千二百四十七条。

\begin{table}[htbp]
\centering
\caption{六卷的名目与条数}
\begin{tabular}{@{}lll@{}}
\toprule
卷 & 内容 & 条数\\
\midrule
居止 & 房舍、通道、环界 & 二百八十一 \\
出入 & 门的开法与值守 & 一百七十四 \\
岁计 & 口粮、库藏、支给 & 二百六十六 \\
器用 & 工具、备件、修缮 & 二百三十 \\
名籍 & 人口、迁转、除籍 & 一百九十二 \\
净秽 & 清除、处置、封存 & 一百零四 \\
\bottomrule
\end{tabular}
\tabsource{条数为重排后的计数。原册的次序已不可知，各卷之间的先后也只是按内容推定。}
\end{table}

六卷之中，篇幅最厚的是居止，最薄的是净秽。这正与常理相左：在一座聚居十余万人的城池里，照理说管人头的名籍应当比管房舍的繁复，管清扫的琐务应当比修缮细碎，可这里的厚薄次序偏偏反着来。

\section{题头的三个字}

第一枚金属片置于卷首，篇幅比其余各枚都短，只有三行，字迹却刻得格外用力。题头三字作小篆体，与秦始皇二十六年诏版上的字体几乎同形，释读作「月之都」；三字下方还补刻了四个小字，「其城自谓」，表明这是居住其间的人对这座城的自称。\footnote{小篆的辨认依据，是秦始皇二十六年诏版上那一路写法：笔画匀圆，结构修长。此三字与诏版同形，释读因此较有把握。}

本书此后即沿用这一称谓。它与前两章用到的几个地名性质迥异：静海七号是现代编目时赋予的编号，「有穷氏旧地」是借用传世文献里的古族名，唯独「月之都」三字出自材料本身，且附有刻手的自注。同一批材料还给出了居民的称呼：名籍卷首条的页眉上刻着「月之民」三字，该处按例本该只题卷名「名籍」二字，刻手却在旁另注一句「其自称也」。城名「月之都」，人名「月之民」，皆非笔者杜撰，而是材料自署的本名；其余诸如静海七号之类，均系外加的代号。

\section{空着的修订栏}

全书一千二百四十七条，每条末尾均留出一栏空白。在寻常的官府文书或律令中，这一栏专用来登记修订记录：某年某人改了何字、因何而改。可这批金属片上的修订栏，自始至终无一字落笔。

全栏皆空，按理有两种解释：一是两万年间条文从未改动，因而无从下笔；二是私下改过而未作记录。前者要求一部管辖十余万人、历时两万年的法典一字不易，后者则要求六卷上下竟无一人肯记上一笔。单凭文字本身，这两种可能实在难以分辨。

能帮上忙的是刀口痕迹。全书一千二百四十七条，共辨出十一类不同的刀口。\footnote{刀口的判别按入刀深度、槽壁倾角与收刀处的残留三项，三人以上各自复核过。同一刻手在不同卷里的写法有细微出入，故数字可能偏多一两个。}十一类刀口说明两万年间金属片曾重刻或补刻过多次，刻手的槽宽与入刀深浅各有不同，可所有条文的字句却始终一字不改。局部几条不改尚可说是恪守祖制，六卷一千二百四十七条在重刻数次之后依然一字未易，只能说明在他们的认知里，没有任何人认为需要改动。至于为何无人动过修改之念，本章尚不能答，须待第四章与第五章讲清他们判断轻重的机制为何在两万年前停摆，方能解开。

\section{缺的两卷}

将六卷名目摊开来看，立刻会发现少了两样极要紧的东西：既没有「医」这一卷，也没有「育」这一卷。一座住着十余万人的城池，历经两万年而无医无育，这等缺口已非疏漏二字可以搪塞。不过单凭目录空缺，尚只能说明章程不管这两件事，还不能直接推出城里的人不会生病、不会生育；要从「不管」推到「没有」，还得看另一件实物证据，即第四百零九枚的背面。

第四百零九枚属于名籍卷。其正面记录一人，姓名一格，定日一格；背面本是供吏员登记职务迁转之用，上面却接连写了六行：定日之后第三年补外工，第四十年调入器用，第一百二十年掌库，第二百二十年调回外环，至第三百二十七年除籍，除籍理由写作「失足」。一个人在同一座城里连续任职三百二十七年，其间既无致仕，也无病退或「以老乞骸骨」一类字样，章程里显然根本未曾给衰老与疾病留出位置。

岁计卷的条文也可作旁证：口粮按「人日二升」支给，两万年未曾变更，既不因年老而减，也不因卧病而增。名籍卷中凡记死亡，一概归入「削」字一栏，理由只有事故、逃亡、除籍三类，同样查不到一条老病。三处材料合观，已足以证明章程中确实没有老病之设，且确有个案在职三百二十七年。诚然，一个人或可偶然长寿，全城却断无此理；但要将全城不老坐实，仍须等到第五章讲清净秽制度之后，论证才算完整。

\section{字、尺与形制}

这批金属片所透露的营建与制度信息，可分字、尺、形制三项来看。先看这三项在月之都内各自呈何面貌，再看它们与地面历史有何牵连。

文字是刻上去的，而非写就。四百一十二枚单面刻字，字口深约零点一毫米，宽约零点四毫米，一笔之内粗细变化不及十分之一，这种均匀度绝非纯凭手腕所能达成，必有一套固定刀法支撑。同一批片子上，有二十一处笔画起止带有一个极小的三角形尖，位置固定且形态一致，如同刻刀留下的签名，足见这二十一处出自同一把或同一种刀具。换言之，写稿者先书，刻手照刻，刻手虽换了几代，刀法却始终未变。

尺度则直接见于条例。移交清单记门宽与环距，用的是「里」与「丈」；居止卷记房舍开间与进深，同样用这两个单位。按十尺为丈、三百步为里的进制，内环直径一点二公里、垣高十二丈，放在这套尺度里丝丝入扣。

形制一半见于居止卷的文字，一半则留在轨道摄影中。自六十年代后期起，月球轨道器与各次登月摄影曾将静海七号拍过许多遍，台基、屋脊与门道皆清晰可辨，底片至今仍保存在公开档案里。受记忆消除场之故，学界从未有一份文献提及这些院落，也无人拿着照片去追问过半句。\footnote{记忆消除场本身不在本书范围，见丛书《记忆消除场的文献学后果》。}从影像上看，静海七号内部是一片片规整的院落：台基高出地面，四面设阶，院内分作几进，屋舍之间以廊相连，屋面成排的坡顶朝向一致，屋脊线笔直如削。月面无雨，坡顶在此全无排水之用，它却依旧照着原来的样式盖了起来。居止卷二百八十一条，一行一条，只记宽若干丈、深若干丈、阶若干、门若干，不加修饰，也不解释缘由，读来宛如一部工书；而卷中所记的台基尺度与台阶数目，正与影像上量出的结果相符。不过这里能对齐的仅限于整体格局，至于举折、出檐等须靠精细分寸比对的细部，在一公里分辨率的影像上还看不真切。

若把字、尺、形制三项并置一处，它们竟与公元前二二一年秦始皇统一天下后在地面推行的那一套完全同路。文字是小篆，与秦始皇二十六年诏版如出一辙；\footnote{秦始皇二十六年诏版，铭文起首「廿六年皇帝尽并兼天下诸侯」，传世与出土者不止一件，字口一致。}尺度按一里四百一十五点八米、一丈二点三一米折算，内环直径与垣高俱通，用的是同一套进制与同一把尺；形制则是高台基、多进院落、坡顶成排，与中原营造法式毫无二致，折算成米亦全都落在秦汉常用尺度区间之内。

三样都对得上，唯独年代差着两万年。按同批器物与岩芯层位，金属片的年代在距今两万一千余年上下；而小篆、这套尺与这种院落格局，在地面上是公元前二二一年之后才定型的。面对这处巨大的时间错位，大抵有三种解释可以一试。

第一种解释是两地各自独立演化，最终碰巧长成同一模样。文字既有约定俗成之分，也有随手取舍之别：一个字如何结体、如何起笔收锋，原本不止一种写法，选定哪一种与字义无关，端赖谁先落笔、后来者依谁临摹。两地若从未往来，各自摸索出几十上百个结体完全一致的小篆字，概率已然微乎其微；再要让尺度的进制各自凑成十尺为丈、三百步为里，让无雨的月面与多雨的中原各自长出同一套高台坡顶院落，三者同时巧合，这条解释显然撑不住。

第二种解释是这批金属片乃后人照着小篆重录的抄本。先不论抄本抄不出尺与城，单看金属片本身的工艺也过不了关：其合金成分与月面样品同类，以今日冶炼手段尚且配制不出；片厚不足一毫米，一千二百四十七条密密麻麻刻于其上而无一开裂；字口深浅误差不超过零点零二毫米，收刀残留形状亦全篇一致。若要重录，必先重造如此薄片，而今日根本造不出来，这一条也可以排除。

于是只剩第三种解释：两地同源，字、尺与形制皆由一处流向另一处。前两说既已被材料排除，这一说尚无材料与之相抵；至于流向何方，片子早于小篆两万年，工艺也远在其上，只能是地面自月面得之。

当然，前两条被排除，与第三条被证实，中间还隔着很远的距离，笔者心里有数。一位英国作家笔下的侦探曾说，排除一切不可能之后，剩下的无论多么荒谬便是真相。\footnote{语出《四签名》。}就本章现有的材料而言，推理的分量也只能走到「未被排除」这一步，至于更扎实的论证，还得留待后续章节去补。

\section{本章的限度}

通盘检视，本章依赖的不过三样东西：四百一十二枚金属片、一千二百四十七条空着的修订栏，以及第四百零九枚的背面。这三样东西全出在同一只来历不明的木箱里。

单一来源的局限是显而易见的：修订栏空白只说明「未曾写下改动」，并不能直接等同于「绝无改动」；第四百零九枚背面也只记了一人一例，孤证难以推及全城。至于小篆与尺度究竟谁先谁后、如何下传，本章只把现象摆在台面上，第七章重逢历法残卷时，才会有一件更硬的证据来接住这个话头。
